\section{Self e a origem do JIT da HotSpot}
\label{sec:hist-self}

O Self era uma linguagem dinâmica, sem classes, em que tudo é envio de
mensagem. Isso deveria torná-la lenta. O grupo de Stanford (e depois da
Sun) passou anos provando que não precisava ser, e as técnicas que
inventou para isso estão na HotSpot.

\paragraph{1989: customização e compilação dinâmica.}
\textcite{paper:self-efficient-impl} agrupam objetos clonados do mesmo
protótipo em \emph{maps} (o que dá ao compilador informação de tipo) e
compilam em runtime várias versões de um mesmo método, cada uma
especializada para o \emph{map} do receptor, o que permite fazer inline
das mensagens enviadas a \texttt{self}. Os autores afirmam que essa
implementação roda duas vezes mais rápido que o Smalltalk mais rápido da
época.

\paragraph{1991: inline caches polimórficos.}
O inline cache da seção anterior falha quando um ponto de chamada vê
vários tipos. \textcite{paper:self-pics} propõem os
\termo{polymorphic inline caches}{PICs: inline caches com uma entrada
por tipo de receptor já visto naquele ponto de chamada.}, e notam o
efeito colateral mais importante: o PIC \emph{registra} os tipos que
apareceram em cada ponto de chamada, e o compilador pode usar isso ao
recompilar. Números dos autores, medidos no próprio sistema Self: mediana
de 11\% de ganho só com PICs e de 27\% quando o compilador usa os tipos
coletados, com cerca de metade das chamadas não inlinadas eliminadas.

\paragraph{1992: desotimização.}
Código muito otimizado é difícil de depurar.
\textcite{paper:self-deoptimization} resolvem isso
\enquote{dynamically deoptimizing code on demand}: só as ativações que
estão sendo depuradas voltam a código não otimizado, e o resto roda em
velocidade máxima. Com isso, também dá para mudar um programa em
execução e ver o efeito na hora. Uma ferramenta de depuração virou a
\termo{deoptimization}{Desotimização: trocar, em plena execução, frames de
código otimizado por frames equivalentes de código não otimizado.} que a
HotSpot usa quando uma especulação falha
(\citetitle{artigo:shipilev-uncommon-traps}).

\paragraph{1994: otimização adaptativa.}
A tese de \textcite{paper:holzle-thesis-adaptive-optimization} junta as
peças: um compilador rápido gera o código inicial, as partes muito usadas
são recompiladas por um compilador otimizador usando o \emph{type
feedback} e o código não otimizado é recriado quando necessário. Leia de
novo a seção~\ref{sec:jit}: interpretador ou C1 primeiro, C2 com perfil
depois, \emph{uncommon trap} e desotimização quando o perfil mente. É o
mesmo desenho.

\paragraph{A ligação, nas palavras de quem fez.}
Não precisamos adivinhar a ligação. O paper do compilador C2,
\textcite{paper:hotspot-server-compiler}, diz que o runtime usa
otimização adaptativa para concentrar a compilação nos métodos críticos
e cita Hölzle e Ungar para isso; e diz que, quando o carregamento de uma
classe invalida uma decisão de inlining, os métodos dependentes são
desotimizados, citando o paper de 1992 acima.
\textcite{paper:truffle-partial-evaluation}, da Oracle Labs, abrem o
paper dizendo que VMs como a HotSpot e o V8 \enquote{follow the design
that was first implemented for the SELF language}.

\paragraph{O que isso explica na JVM.}
\begin{itemize}
\item Por que o warm-up existe e por que a HotSpot coleta perfil: o
  compilador otimizador só é bom porque viu os tipos reais
  (seção~\ref{sec:jit}).
\item Por que chamadas virtuais e interfaces não são caras na prática:
  inline caches e inlining guiado por tipo, herdados do Self.
\item Por que carregar uma classe nova pode deixar o programa mais lento
  por um instante: desotimização de métodos que assumiam um só
  implementador (seção~\ref{sec:class-loading}).
\item Por que o nativo pode perder em throughput de pico, como na
  medição do Quarkus (seção~\ref{sec:jvm-vs-nativo}), embora o paper do
  Native Image fale em preservar o desempenho de pico: sem perfil da
  execução real e sem desotimização, o compilador só pode apostar no que
  prova no build (seção~\ref{sec:vms-jit-aot}).
\end{itemize}
